In the canonical coordinates of the Hamilton-cycle side, the following
three cuts give four additional fixed-$K$ interfaces. Coordinates are
zero-based, and $S$ denotes the actual component vertices satisfying
the displayed predicate.
\[
\begin{array}{c|c|c|c}
B&S&|\delta_B(S)\cap((K\cap E(B))\setminus P_B)|
  &\delta_H(S)\\\hline
00121&x_0x_1=23,\ x_3\ne0&12&23203\text{--}23213\\
01024&x_1x_2x_3=321&10&10212\text{--}13212\\
01024&x_0x_1=23&15&13130\text{--}23130
\end{array}
\]
Here $H=B-((K\cap E(B))\setminus P_B)$. In the first row,
$K$ is the printed $00101$ protected set transported by the original
map $(0,+1,0132)$ of its edit to colour $3$ at coordinate $3$.
In the second row it is the printed $01014$ protected set transported
by $(3,-1,2031)$ of its edit to colour $3$ at coordinate $2$.
In the last row it is the printed $01234$ protected set transported
by $(0,+1,3102)$. This is both the forward $01234$ edit to colour $2$
at coordinate $0$ and the reverse role of the $01024$ edit to colour
$3$ there, whose original map $(0,+1,2130)$ is the inverse.
All transported sets retain their actual reserve-pool conditions.
Each nontrivial cut in $H$ has only one edge, so $H$ has no spanning
cycle. This does not require the remaining bridge itself to avoid
the reserve pool. The $12,10,15$ cross-edge memberships and singleton
remainders are actual finite checks, not a new uniform classification.
